\section{Giá trị của liên minh không tự có}
\label{sec:ch04-gia-tri-lien-minh-va-phan-phoi-nen}

\index{phân phối nền}%
\index{đặc trưng phụ thuộc}%
Mục trước đã cố định cách chia khi $v(S)$ đã biết. Với một mô hình dự đoán,
việc khó hơn là gán nghĩa cho liên minh $S$: đầu vào chỉ giữ các đặc trưng trong
$S$ thì mô hình sẽ trả gì? Hầu hết mô hình không nhận một đầu vào bị khuyết
tùy ý, nên bài báo SHAP gốc thay giá trị đó bằng kỳ vọng có điều kiện của đầu
ra mô hình khi biết các đặc trưng trong $S$~\autocite{p10shap}.

Nói cách khác, mỗi liên minh được đánh giá trên một \tn{phân phối nền}{background
distribution}. Phân phối nền trả lời câu hỏi những đặc trưng còn thiếu nên trông
như thế nào. Nếu giữ tuổi và bỏ thu nhập trong một hồ sơ khách hàng, ta có thể
lấy trung bình trên các mức thu nhập từng đi cùng độ tuổi đó. Đây là một mệnh
đề về dữ liệu, không phải hệ quả của các tiên đề Shapley.

Bài báo gốc cũng nêu một xấp xỉ rẻ hơn: thay kỳ vọng có điều kiện bằng kỳ vọng
biên, tức lấy các đặc trưng còn thiếu từ phân phối biên của chúng thay vì từ
phân phối có điều kiện trên phần đã biết.
Xấp xỉ này có thể tạo các tổ hợp mà dữ liệu thật hiếm khi hoặc không bao giờ
có. Vì vậy, khi hai đặc trưng phụ thuộc nhau, một giá trị SHAP không tự cho biết
liệu nó đang mô tả quan hệ dữ liệu quan sát được hay một can thiệp giả định.

Các công trình sau đó phân biệt rõ cách nhìn có điều kiện và cách nhìn biên.
Trong cách nhìn biên, các đặc trưng còn thiếu được lấy từ phân phối biên; bài báo
về suy luận thống kê cho SHAP lưu ý rằng các phần mềm thường dùng cách này và
rằng hai cách nhìn có những diễn giải khác nhau~\autocite{p19shapinference}.
Vì vậy, một biểu đồ SHAP cần đi cùng câu trả lời về phân phối nền và về
\tn{quy tắc hoàn thiện}{completion rule}, tức cách điền giá trị cho các đặc
trưng còn thiếu của một liên minh trước khi đưa liên minh ấy vào $f$, chứ
không chỉ danh sách các $\phi_i$. Hai thứ ấy tách nhau: quy tắc hoàn thiện nói
giá trị thay thế được lấy theo phân phối có điều kiện, theo phân phối biên, hay
bằng một giá trị cố định, còn phân phối nền là tập dữ liệu mà quy tắc ấy lấy
mẫu từ đó.
